\chapter{Literature Review}
\label{ch:lit}

\section{Honeypots}

Spitzner defines a honeypot as ``an information system resource whose value lies in
unauthorised or illicit use of that resource''~\cite{spitzner2002honeypots}. Because a honeypot
has no legitimate users, every interaction with it is suspect, which gives it a very low
false-positive rate compared with detection systems deployed on production traffic.

Honeypots are usually classified by the degree of interaction they allow:

\begin{description}
  \item[Low interaction] systems emulate a service only up to the point of recording a
    connection or a login attempt. Honeyd~\cite{provos2004honeyd} emulates whole networks of
    such hosts. They are cheap and safe but reveal little about post-compromise behaviour.
  \item[High interaction] systems give the attacker a real operating system, usually a
    virtual machine that is monitored and reset. They capture the most realistic behaviour, but
    they are expensive to operate and risky, because the attacker controls a real machine that
    can be used to attack others.
  \item[Medium interaction] systems emulate an operating system's user-facing behaviour (a
    shell, a filesystem, command output) in software. Attackers can run commands, but nothing
    they type is executed. This is the design point of Kippo, Cowrie~\cite{cowrie} and of this
    project.
\end{description}

\section{SSH Honeypots}

Kippo, and its maintained successor Cowrie~\cite{cowrie}, are the most widely deployed
medium-interaction SSH honeypots. Cowrie emulates a Debian system with a fake filesystem,
logs login attempts and commands, records sessions for replay, and stores files that attackers
download. Its emulation is \emph{static}: responses depend only on the command and the state of
the fake filesystem. Studies built on such honeypots consistently find that most SSH
traffic is automated: a few credential pairs dominate, and post-login activity is a short
script of reconnaissance followed by a download-and-execute step. Interactive human sessions
are rarer but longer and more varied. Our attacker model (Chapter~\ref{ch:formulation})
reflects these qualitative findings.

\section{Adaptive Honeypots}

The idea that a honeypot should \emph{choose} its responses predates the use of RL. Wagener
et al.'s Heliza~\cite{wagener2011heliza} is a high-interaction honeypot that can allow, block,
substitute or insult attacker commands. It models the interaction as a game and uses
reinforcement learning to pursue goals such as collecting attacker tools or wasting their time.

Pauna and Bica's RASSH~\cite{pauna2014rassh} brought the idea to a medium-interaction SSH
honeypot built on Kippo. Its agent chooses among allowing, delaying, faking and insulting
commands, with SARSA learning a policy that rewards prolonged interaction. A later version
replaced the table with a deep Q-network~\cite{pauna2018qrassh}. Dowling, Schukat and
Barrett~\cite{dowling2018conceal} used RL to keep attackers engaged with an IoT honeypot and
to conceal honeypot functionality.

Two observations from this literature shape the present work.

\begin{enumerate}
  \item \textbf{The action sets are similar.} Allow, delay, fake and block recur. ``Insult''
    also appears, but it gives the honeypot away, so we replace it with a \emph{tarpit}, which
    dribbles output slowly and is an established deception technique against automated tools.
  \item \textbf{Evaluation is the weak point.} Adaptive honeypots are typically evaluated by
    deploying them for some weeks and reporting engagement statistics. That shows the system
    works, but it cannot compare policies fairly (different policies meet different attackers
    on different days). It also cannot say how far a learned policy is from what is achievable.
    Our contribution is an evaluation methodology that does both, using exact computation on a
    generative attacker model.
\end{enumerate}

\section{Reinforcement Learning Background}

Reinforcement learning~\cite{sutton2018rl} studies agents that learn to act from rewards
through interaction. When the environment is a finite Markov decision process with unknown
dynamics, \emph{Q-learning}~\cite{watkins1992q} learns the optimal action-value function from
samples using the temporal-difference update
\[
Q(s,a) \leftarrow Q(s,a) + \alpha\big[r + \gamma \max_{a'} Q(s',a') - Q(s,a)\big].
\]
Watkins and Dayan proved convergence to the optimum with probability one, provided every
state--action pair is updated infinitely often and the step sizes satisfy the Robbins--Monro
conditions $\sum_n \alpha_n = \infty$ and $\sum_n \alpha_n^2 < \infty$~\cite{robbins1951}.
Even-Dar and Mansour~\cite{evendar2003rates} showed that the \emph{rate} of convergence depends
strongly on the schedule: polynomial rates $\alpha_n = n^{-\omega}$ with $\tfrac12 < \omega < 1$
converge polynomially, whereas $\alpha_n = 1/n$ can need exponentially many samples. In
practice a constant step size is common, because it keeps tracking a changing environment. The
price is that the estimate never settles. Chapter~\ref{ch:results} shows that in our setting
this price is decisive.

\emph{SARSA}~\cite{sutton2018rl} is the on-policy counterpart of Q-learning: it bootstraps from
the action actually taken next rather than the greedy one. \emph{Double
Q-learning}~\cite{vanhasselt2010double} removes the upward bias introduced by the max operator.

When the agent observes only part of the state, the problem becomes a partially observable
MDP (POMDP). Singh, Jaakkola and Jordan~\cite{singh1994pomdp} showed that the best policy over
observations alone may need to be stochastic, and that TD methods applied directly to
observations may converge to sub-optimal policies or not at all. A honeypot cannot see the
attacker's type or intent, only their commands, so this result is directly relevant.

\section{Event Streaming and Storage}

Separating event capture from processing through a durable log is a standard architecture
for telemetry systems. Redis Streams~\cite{redisstreams} provide an append-only log with
consumer groups, per-consumer pending-entry lists and explicit acknowledgement. This supports
at-least-once delivery without a heavier platform such as Kafka. PostgreSQL provides the
relational model, transactions and indexing that investigative queries need.

\section{Research Gap}

Existing adaptive SSH honeypots demonstrate feasibility but report neither how their learned
policy compares with the best achievable policy nor how the benefit depends on the attacker
population. This project contributes (i) a formal model in which those quantities are
defined, (ii) an exact method to compute them, and (iii) an honest measurement of the benefit,
including the conditions under which a static honeypot is already near-optimal.
